Megapascal (MPa) Ranges for Human Aortic Wall Tensile Strength

Ultimate tensile strength (UTS, also called failure stress) of the human aortic wall is anisotropic (stronger circumferentially than longitudinally) and varies substantially with age, disease state, anatomic location, testing direction, and loading rate. Values are typically reported in kilopascals (kPa) or megapascals (1 MPa = 1000 kPa). The table below synthesizes representative ranges from biomechanical studies of fresh or banked human tissue (primarily ascending and descending thoracic aorta, with some abdominal data).

| Condition / Tissue Type                  | Direction          | Typical UTS Range (MPa) | Notes / Key Observations |
|------------------------------------------|--------------------|--------------------------|--------------------------|
| Healthy / Non-aneurysmal (younger adults) | Circumferential   | 1.2 – 2.2 (up to ~3–5 under dynamic loading) | Highest strength; Windkessel-capable tissue |
|                                          | Longitudinal      | 0.6 – 1.7               | Consistently lower than circumferential |
| Healthy / Non-aneurysmal (older adults) | Circumferential   | 1.0 – 1.8               | Age-related stiffening + modest strength loss |
|                                          | Longitudinal      | 0.5 – 1.2               | Greater relative decline with age |
| Aneurysmal / Atherosclerotic         | Circumferential   | 0.5 – 1.7 (often 0.8 – 1.5) | Reduced vs. age-matched healthy; high inter-specimen variability |
|                                          | Longitudinal      | 0.3 – 1.0               | Frequently the weaker direction for initiation of transverse tears |
| Dissected media / Severely degraded  | Circumferential   | 0.4 – 1.3               | Markedly lower; delamination strength often << 1 MPa |
|                                          | Longitudinal      | 0.2 – 0.8               | Lowest values; aligns with vulnerability thresholds |
| Layer-specific (media focus)         | Circumferential   | ~1.2 – 1.5 (healthy media) | Media is primary load-bearing layer for dissection resistance |
|                                          | Longitudinal      | ~0.5 – 0.7              | Media failure often initiates longitudinally under physiologic biaxial stress |
| Radial (through-wall / peel)         | Radial            | 0.1 – 0.3               | Much weaker; relevant to delamination / false-lumen propagation |

Key Context for the Balloon-Expansion / Jet-Hammer Thesis
Healthy reference**: Circumferential UTS commonly 1.5–2.5 MPa (occasionally higher under rapid loading). Longitudinal values are typically 40–60 % lower.
Degraded / aging / clotted wall: Mean UTS frequently falls into the 0.5–1.2 MPa range, with many specimens (especially longitudinal orientation or transition zones at plaque edges) dropping below **0.8 MPa. This matches the thesis threshold at which cyclic ballooning stress + high-velocity jet impact can exceed tissue strength.
Physiologic wall stress** under normal blood pressure is far lower (typically 0.1–0.3 MPa), providing a large safety factor in healthy tissue. In stiffened, matrix-degraded segments the safety factor shrinks dramatically, so modest pressure surges or localized jet forces become sufficient to initiate an intimal micro-tear.

These ranges are derived from uniaxial and biaxial tensile testing of human specimens and show substantial scatter; individual patient tissue can deviate outside the tabulated intervals. The progressive drop from healthy multi-megapascal strength to sub-megapascal values in enzymatically and mechanically preconditioned aortas is the central material change that enables the balloon-expansion and high-velocity jet mechanisms to produce acute dissection.
